% ==============================
% Chapter 2 – Problem Definition
% ==============================

\chapter{Problem Definition}
% --------------------------------
% 2.1 Problem Statement
% --------------------------------

\section{Problem Statement}

\textbf{Writing Guidelines (100–150 words)}

Clearly define the core problem that the proposed system aims to solve. 
Explain existing challenges, inefficiencies, limitations, or gaps in the 
current system or process. Identify affected stakeholders such as users, 
organizations, or administrators. Highlight why this problem is significant 
in technical, operational, or societal terms. Avoid describing the proposed 
solution here — focus strictly on defining the problem in a clear, logical, 
and evidence-based manner.

\vspace{0.5cm}

\textbf{\textit{Replace this text with your finalized Problem Statement before submission.}}



% --------------------------------
% 2.2 Proposed Solution Overview
% --------------------------------

\section{Proposed Solution Overview}

\textbf{Writing Guidelines (100–150 words)}

Provide a high-level conceptual description of the proposed system. 
Explain how it addresses the issues identified in the Problem Statement. 
Mention major features, innovations, or technologies used. Describe 
the value proposition and expected improvements in efficiency, usability, 
performance, automation, or transparency. Keep the explanation conceptual 
and avoid detailed implementation or technical architecture at this stage.

\vspace{0.5cm}

\textbf{\textit{Insert the overview of your proposed solution here.}}



% --------------------------------
% 2.3 Objectives of the Proposed System
% --------------------------------

\section{Objectives of the Proposed System}

\textbf{Writing Guidelines}

List clear, measurable, and verifiable objectives. Each objective must:
\begin{itemize}
    \item Begin with an action verb (e.g., Develop, Provide, Improve, Ensure).
    \item Be specific and testable.
    \item Directly align with the Problem Statement.
    \item Be achievable within project constraints.
\end{itemize}

\vspace{0.3cm}

\textbf{Business Objectives (BO):}

\begin{itemize}
    \item BO-1: Develop a system that ...
    \item BO-2: Improve ...
    \item BO-3: Provide ...
    \item BO-4: Ensure ...
    \item BO-5: Automate ...
\end{itemize}



% --------------------------------
% 2.4 Scope of the System
% --------------------------------

\section{Scope of the System}

\textbf{Writing Guidelines (100–150 words)}

Define the boundaries of the project clearly. Specify what the system 
will include and what it will not include. Identify the target users, 
supported platforms (web, mobile, desktop), and core functionalities. 
Clarify operational limits and project coverage. Distinguish scope 
from limitations — scope defines coverage, while limitations define 
constraints.

\vspace{0.5cm}

\textbf{\textit{Insert scope description here.}}



% --------------------------------
% 2.5 System Architecture and Modules
% --------------------------------

\section{System Architecture and Modules}

\textbf{Writing Guidelines (100–150 words)}

Provide an overview of the system architecture. Specify whether the 
system follows client-server, three-tier, MVC, microservices, or 
another architectural pattern. Explain why modular design is used. 
Describe how modules interact and exchange data. A high-level system 
architecture diagram should be included in this section.

\vspace{0.5cm}

\textbf{\textit{Insert architecture overview here.}}



% --------------------------------
% 2.5.1 Module Descriptions
% --------------------------------

\subsection{Module 1: Profile Management}

\textbf{Description (80–120 words)}

Explain the purpose of this module. Describe user registration, 
authentication, authorization, and profile management features. 
Specify supported user roles.

\textbf{Functional Requirements:}

\begin{itemize}
    \item FE-1: The system shall allow users to register using valid credentials.
    \item FE-2: The system shall authenticate users before granting access.
\end{itemize}



\subsection{Module 2: Communication Module}

\textbf{Description (80–120 words)}

Describe how users interact within the system. Mention messaging, 
notifications, scheduling, or real-time communication capabilities.

\textbf{Functional Requirements:}

\begin{itemize}
    \item FE-1: The system shall enable secure messaging between users.
    \item FE-2: The system shall generate real-time notifications.
    \item FE-3: The system shall maintain communication logs.
\end{itemize}



\subsection{Module 3: Core Functional Module}

\textbf{Description (80–120 words)}

Explain the primary operational features that deliver the core 
value of the system. This module should directly address the 
problem identified in Section 2.1.

\textbf{Functional Requirements:}

\begin{itemize}
    \item FE-1: The system shall process user requests efficiently.
    \item FE-2: The system shall store and retrieve data accurately.
    \item FE-3: The system shall generate relevant outputs based on user input.
\end{itemize}



\subsection{Module 4: Management / Admin Module}

\textbf{Description (80–120 words)}

Explain administrative controls, monitoring dashboards, reporting 
tools, and configuration management features.

\textbf{Functional Requirements:}

\begin{itemize}
    \item FE-1: The system shall allow administrators to manage user accounts.
    \item FE-2: The system shall generate system usage reports.
\end{itemize}



\subsection{Module 5: Advanced / Additional Features}

\textbf{Description (80–120 words)}

Describe advanced capabilities such as AI-based recommendations, 
data analytics, automation, payment gateway integration, or external 
API integration.

\textbf{Functional Requirements:}

\begin{itemize}
    \item FE-1: The system shall integrate with external APIs securely.
    \item FE-2: The system shall provide analytics-based insights.
\end{itemize}


% --------------------------------
% 2.7 Assumptions and Dependencies
% --------------------------------

\section{Assumptions and Dependencies}

\textbf{Writing Guidelines (80–120 words)}

List assumptions made during system design. Identify dependencies 
on external systems, APIs, hardware, or third-party services.

\begin{itemize}
    \item The system assumes users have stable internet connectivity.
    \item The system depends on third-party API availability.
    \item The system requires compatible hardware/software environments.
\end{itemize}

%%%%%%%%%%%%%%%%%%%%%%%%%%%%%%%%%%%%%%%%%%%%%%%%%%%%%%
\section{Chapter Summary}
%%%%%%%%%%%%%%%%%%%%%%%%%%%%%%%%%%%%%%%%%%%%%%%%%%%%%%

This chapter defined the core problem addressed by the proposed system and 
outlined the conceptual solution approach. Clear objectives were established 
to ensure alignment with identified challenges. The system scope and 
architectural overview were presented to define structural boundaries 
and major modules. Key assumptions and dependencies were also documented 
to clarify operational constraints. 

This chapter provides the problem foundation and sets the direction for 
the detailed software design presented in the next chapter.
